\documentclass[12pt,a4paper]{article}
\usepackage{ctex}
\usepackage{amsmath,amscd,amsbsy,amssymb,latexsym,url,bm,amsthm}
\usepackage{epsfig,graphicx,subfigure}
\usepackage{enumitem,balance}
\usepackage{wrapfig}
\usepackage{mathrsfs,euscript}
\usepackage[usenames]{xcolor}
\usepackage{hyperref}
\usepackage[vlined,ruled,commentsnumbered,linesnumbered]{algorithm2e}

\newtheorem{theorem}{Theorem}
\newtheorem{lemma}[theorem]{Lemma}
\newtheorem{proposition}[theorem]{Proposition}
\newtheorem{corollary}[theorem]{Corollary}
\newtheorem{exercise}{Exercise}
\newtheorem*{solution}{Solution}
\newtheorem{definition}{Definition}
\theoremstyle{definition}


%\numberwithin{equation}{section}
%\numberwithin{figure}{section}

\renewcommand{\thefootnote}{\fnsymbol{footnote}}

\newcommand{\postscript}[2]
 {\setlength{\epsfxsize}{#2\hsize}
  \centerline{\epsfbox{#1}}}

\renewcommand{\baselinestretch}{1.0}

\setlength{\oddsidemargin}{-0.365in}
\setlength{\evensidemargin}{-0.365in}
\setlength{\topmargin}{-0.3in}
\setlength{\headheight}{0in}
\setlength{\headsep}{0in}
\setlength{\textheight}{10.1in}
\setlength{\textwidth}{7in}
\makeatletter \renewenvironment{proof}[1][Proof] {\par\pushQED{\qed}\normalfont\topsep6\p@\@plus6\p@\relax\trivlist\item[\hskip\labelsep\bfseries#1\@addpunct{.}]\ignorespaces}{\popQED\endtrivlist\@endpefalse} \makeatother
\makeatletter
\renewenvironment{solution}[1][Solution] {\par\pushQED{\qed}\normalfont\topsep6\p@\@plus6\p@\relax\trivlist\item[\hskip\labelsep\bfseries#1\@addpunct{.}]\ignorespaces}{\popQED\endtrivlist\@endpefalse} \makeatother



\begin{document}
\noindent

%========================================================================
\noindent\framebox[\linewidth]{\shortstack[c]{
\Large{\textbf{Exercise Sheet 12}}\vspace{1mm}\\
Discrete Mathematics by Hongfei Fu, 2023.11.20}}

\begin{center}
	\footnotesize{\color{black} \quad Name:\_\_\_\_\_\_\_\_\_  \quad Student ID:\_\_\_\_\_\_\_\_\_ \quad Class: \_\_\_\_\_\_\_\_\_\_\_\_} \\


\end{center}

\begin{enumerate}

\item ([R], Page 676-677, Exercise 36, 37, 42, 43(\textbf{2pt}), 44(\textbf{2pt})) In Exercises 36, 37, 42, 43, 44 determine whether the given pair of graphs is isomorphic. Exhibit an isomorphism or provide a rigorous argument that none exists.\\

(36). \includegraphics*[scale=0.7]{36@p676.png} \\
\textbf{Answer Area:}\\
No. A proof is as follows:
$\because$ In the left graph, there isn't a vertex whose degree is 4. But in the right graph, deg($v_{2}$) = 4.\\
$\therefore$ The two graphs are not isomorphic.
\vspace{1cm}

(37). \includegraphics*[scale=0.7]{37@p676.png} \\
\textbf{Answer Area:}\\
Yes. There is a bijection:\\
$f(u_{1}) = v_{1}$\\
$f(u_{2}) = v_{3}$\\
$f(u_{3}) = v_{5}$\\
$f(u_{4}) = v_{7}$\\
$f(u_{5}) = v_{2}$\\
$f(u_{6}) = v_{4}$\\
$f(u_{7}) = v_{6}$\\
Hence, the two graphs are isomorphic.
\vspace{1cm}

(42). \includegraphics*[scale=0.7]{42@p677.png} \\
\textbf{Answer Area:}\\
No. A proof is as follows:
Assuming that there is a bijection that satisfies the needs.\\
$\because$If $f(u_{i}) = v_{j}$, then $deg(u_{i}) = deg(v_{j})$\\
$\therefore f(u_{5}) = v_{5}$ or $f(u_{5}) = v_{1}.$\\
$\because deg(v_{2}) = deg(v_{4}) = 4 \neg 3 = deg(u_{4}).$\\
$\therefore f(u_{2}) \neg v_{4}$ and $f(u_{2}) \neg v_{2}$.\\
$\therefore f(u_{5}) \neg v_{5}$ and $f(u_{5}) \neg v_{1}$\\
Contradict. The two graphs are not isomorphic.
\vspace{1cm}

(43)(\textbf{2pt}). \includegraphics*[scale=0.7]{43@p677.png} \\
\textbf{Answer Area:}\\
Yes. A bijectionYes. There is a bijection:\\
$f(u_{1}) = v_{1}$\\
$f(u_{2}) = v_{2}$\\
$f(u_{3}) = v_{3}$\\
$f(u_{4}) = v_{7}$\\
$f(u_{5}) = v_{6}$\\
$f(u_{6}) = v_{5}$\\
$f(u_{7}) = v_{10}$\\
$f(u_{8}) = v_{4}$\\
$f(u_{9}) = v_{8}$\\
$f(u_{10}) = v_{9}$\\
Hence, the two graphs are isomorphic.
\vspace{1cm}

(44)(\textbf{2pt}). \includegraphics*[scale=0.7]{44@p677.png} \\
\textbf{Answer Area:}\\

\vspace{4cm}

\end{enumerate}
%========================================================================
\end{document}

